\section*{Proforma}
\begin{tabular}{ll}
  Candidate Number   & 9615Q                                                      \\
  Title              & A type-preserving compiler from a systems language to DTAL \\
  Examination        & Computer Science Tripos -- Part II, 2026                   \\
  Word Count         & 12,992\textsuperscript{1}
  \\
  Code Line Count    & 34,581\textsuperscript{2}
  \\
  Ethics Approval    & Not requested
  \\
  Project Originator & The candidate                                              \\
  Project Supervisor & Yulong Huang                                               \\
  Marking Supervisor & Jeremy Yallop
\end{tabular}

\smallskip
\noindent\textsuperscript{1}Computed with \texttt{texcount}. \textsuperscript{2}Non-test Rust lines under \texttt{veritas/src}.

\subsection*{Original aims of the project}

The project set out to develop a type-preserving compilation toolchain for a systems programming language with refinement types. The goal was to demonstrate that structural safety properties, such as spatial memory safety and control-flow integrity, could be mathematically guaranteed in low-level executable code without trusting the complex compiler itself. This involved designing a source language, building a compiler that lowers high-level constructs to a Dependently Typed Assembly Language (DTAL) while preserving type information, and creating an independent verifier capable of confirming the safety of the resulting DTAL independently of the compiler.

\subsection*{Work completed}

I implemented Veritas: a refinement-typed imperative source language, parser, bidirectional type checker, typed SSA intermediate representation, DTAL backend, independent verifier, physical DTAL pipeline after register allocation, direct x86-64 encoder, runtime intrinsics, and hosted and bare-metal execution targets. The implementation supports contracts, loop invariants, bounded quantifiers, array-bounds checking, borrowing/ownership operations, executable ELF output, and a benchmark and test suite used to evaluate correctness, performance, and trusted-computing-base reduction.

\subsection*{Special difficulties}
None.
